\chapter{Value Iteration and Policy Iteration}

\lecture{Feb. 2}

Bellman Optimality Equation
\[
    V^*(s) = \max_{a} \left( r(s, a) + \gamma \sum_{s'} P(s' | s,a) V^*(s') \right)
\]
The statement that this is a fixed-point has nothing to do with contraction. We want to perform an iterative loop using the Bellman operator. The optimal value function $V^*$ is the unique fixed point of the Bellman operator $T$, i.e., $TV^* = V^*$. Value iteration, defined as $V_{k+1}=TV_k$, with an arbitrary initial value $V_0$ will converge to the optimal value function $V^*$, i.e. $\lim_{k\to\infty} V_k = V^*$.  \\

Optimal policy $\pi^*$ can be computed as
\[
    \pi^*(s) = \argmax_{a} (r(s,a) + \gamma \E_{s'\sim P(\cdot | s,a)}[V^*(s')])
\]
Optimal Q-value function of a policy $\pi$
\[
    Q_\pi(s,a) = \E [\sum_{t=0}^{\infty} \gamma^t r(s_t, a_t) | s_0 = s, a_0 = 0]
\]
Just like optimal value function, you can define an optimal Q-value function
\[
    Q^* (s, a) = \max_{\pi} Q_\pi (s,a)
\]
There exists a conversion from $V^*$ to $Q^*$ through
\[
    Q^*(s,a) = r(s,a) = \gamma \E_{s'\sim P(\cdot | s,a)} [V^*(s')]
\]
And from $Q^*$ to $V^*$
\[
    V^*(s) = \max_{a}Q^* (s, a)
\]
Also from $\pi^*$ to $Q^*$
\[
    \pi^*(s) = \argmax_{a} Q^*(s, a)
\]
In some problems, you may not know the transition probability. \\

Optimal Q-value function $Q^*$ satisfies the Bellman equation
\[
    Q^*(s, a) = r(s,a) + \gamma \E_{s'\sim P(\cdot| s,a)} [\max_{b} Q^*(s',b)]
\]
Define the mapping $F: \R^{|\mathcal{S}||\mathcal{A}|} \to \R^{|\mathcal{S}||\mathcal{A}|}$ as
\[
    (FQ)(s, a) = r(s,a) = \gamma \E_{s'\sim P(\cdot | s,a)} [\max_{b} Q(s', b)]
\]
\section{Policy Iteration}

Recap of policy iteration: at every point, we evaluate how good our situation is (policy, $\pi_k$). Our value function $V_\pi$ satisfies the equation $V_\pi = R_\pi + \gamma P_\pi V_\pi$. How do we find $V_\pi$? \\

Can we find the optimal policy by 'cleverly' searching over the policy space? \\

\noindent Policy iteration: for $k=0,1,2,\dots$ Policy evaluation: compute the value of the policy to get $V_{\pi_k}$
\[
    \text{Solve } T_{\pi_k} V_{\pi_k} = V_{\pi_k}
\]
Policy improvement: perform greedy update to get the next policy $\pi_{k+1}$:
\[
    \pi_{k+1}(x) = \argmax_{a} (r(s,a) + \gE [V_{\pi_k}(s')])
\]
Proposition (convergence of policy iteration): $V_{\pi_{k+a}} \geq V_{\pi_k}$ and $\norm{V_{\pi_{k+1}} - V^*}_{\infty} \leq \gamma \norm{V_{\pi_k} - V^*}_{\infty}$. \\

\noindent Policy iteration:
Lemma (monotone property of $T_\pi$): let $V_1 \geq V_2$. Then, $T_\pi V_1 \geq T_\pi V_2$. \\

\noindent Proof: note that $V_{\pi_k} = T_{\pi_k} V_{\pi_k} \leq TV_{\pi_k} = T_{\pi_{k+1}} V_{\pi_k} $

% TODO: finish all of this
Using the monotone property of $T_\pi$,
\begin{align*}
    V_{\pi_k} \leq T_{\pi_{k+1}} V_{\pi_k} \leq T_{\pi_{k+1}}^{(1)} V_{\pi_k} \leq \dots \leq T_{\pi_{k+1}}^{(2)} \dots
\end{align*}